\section[From Wave-Particle Duality to Quantum States]{From Wave-Particle Duality\\ to Quantum States}\label{sec:from-wave-particle-duality-to-quantum-states}

\subsection{Experimental Motivation: Wave-Particle Duality}\label{subsec:from-wave-particle-duality-to-quantum-states--experimental-motivation-wave-particle-duality}

The purpose of this section is not to give a history of the development of quantum mechanics, but rather to discover \hl{what kind of mathematical theory we will eventually need} to describe the phenomena that we will discuss in these notes.

\highspace
We begin from a very basic question: ``\emph{\textbf{what is light?}}''. Classically, there are two very different possibilities:
\begin{enumerate}
    \item A \definition{Particle} is something localized. At a given instant it \textbf{occupies a particular position}, \textbf{follows a trajectory}, and \textbf{arrives somewhere} as an individual object.

    \item A \definition{Wave}, on the other hand, is \textbf{spatially extended}. Different portion of the wave can overlap and interfere.
\end{enumerate}
Historically, both descriptions were proposed for light. But for our purposes, the important question is not who proposed what. The important question is: \emph{\hl{how can an experiment distinguish wave behaviour from particle behaviour?}} In other words, definitions alone cannot choose between the two descriptions. So we \hl{need an \textbf{experimental criterion} that allows us to distinguish between the two behaviors.}

\highspace
\begin{flushleft}
    \textcolor{Green3}{\faIcon{balance-scale} \textbf{What phenomenon is characteristic of waves? The wave vs. particle descriptions of light}}
\end{flushleft}
Suppose we have light traveling from a source toward some obstacle and then onto a screen. \hl{If light behaved only like classical particles}, the simplest picture would be:
\begin{equation*}
    \text{source} \to \text{particle trajectories} \to \text{screen}
\end{equation*}
\hl{Each particle follows some path, reaches a particular location, and contributes there.}

\highspace
But \textbf{waves behave differently} because different contributions can be added together. For two waves:
\begin{equation*}
    \psi_1, \quad \psi_2
\end{equation*}
The \textbf{combined disturbance}\footnote{%
    The \textbf{total/resultant wave} produced when the two waves are present at the same place and time.
} is:
\begin{equation*}
    \psi_{\text{tot}} = \psi_1 + \psi_2
\end{equation*}
This simple operation has a crucial consequence: the \hl{two waves may reinforce each other or cancel each other.} That phenomenon is called \definition{Interference}.

\newpage

\noindent
So already we have the first important distinction between waves and particles:
\begin{itemize}
    \item \textbf{Particles} are localized objects.
    \item \textbf{Waves} are superposable\footnote{%
            They can be added together. More precisely, in a linear wave system, if $\psi_1$ and $\psi_2$ are allowed waves, then $\psi_1 + \psi_2$ is also an allowed wave.
        } and extended disturbances.
\end{itemize}
At this point, ``superposition'' is still the familiar classical-wave idea.\footnote{%
    The sum of two waves is still a wave.
} We are not yet stating the quantum superposition principle. Also, we have seen that the \textbf{interference phenomenon} is a characteristic of waves, but we have not yet seen any experimental evidence for it.

\highspace
\begin{flushleft}
    \textcolor{Green3}{\faIcon{book} \textbf{Diffraction and interference as a signatures of waves}}
\end{flushleft}
Interference is often observed experimentally through a \definition{Diffraction} pattern. It is one of the \textbf{most convenient experimental manifestations of interference}. In general, interference is the general wave phenomenon, and diffraction is a \hl{particular wave phenomenon in which a wave spreads after an aperture/obstacle, and the resulting diffraction pattern is produced by interference.} So we can have interference without diffraction, for example two independent water waves meeting each other.

\begin{deepeningbox}[: What does a ``wave'' actually have?]
    Before talking about light, think about something more intuitive: a \textbf{wave on a rope}. Imagine a rope that is fixed or held at the ends. If we move one point up and down, the disturbance propagates along the rope. The rope itself does \textbf{not} travel along with the wave. Each small piece of rope mainly moves up and down around its equilibrium position, while the \textbf{wave pattern} propagates horizontally.

    \highspace
    Mathematically, at each point $x$ along the rope, it has some displacement from equilibrium:
    \begin{equation*}
        y(x,t)
    \end{equation*}
    Which is the description of the wave, where $y$ is the vertical displacement of the rope at position $x$ and time $t$ ($y=0$ is the equilibrium position). At one instant maybe $y = 2$ cm, and somewhere else $y = -1$ cm.

    \highspace
    Thus, the quantity $y$ is the \textbf{wave amplitude} at that point and instant. It indicates how far the rope is displaced from its equilibrium position.

    \highspace
    For the case of \textbf{light}, instead of writing $y(x,t)$, we can think about:
    \begin{equation*}
        E(x, t)
    \end{equation*}
    So the ``amplitude'' is the \textbf{value of the oscillating wave quantity}.
\end{deepeningbox}

\highspace
\begin{deepeningbox}[: What happens when two waves meet?]
    This is the defining ingredient behind interference. Suppose two waves arrive at the same place. Their amplitudes add:
    \begin{equation*}
        E = E_1 + E_2
    \end{equation*}
    This is the superposition principle for classical waves. Now imagine at instant:
    \begin{equation*}
        E_1 = +1, \quad E_2 = +1
    \end{equation*}
    Then:
    \begin{equation*}
        E = +2
    \end{equation*}
    They reinforce each other. But perhaps instead:
    \begin{equation*}
        E_1 = +1, \quad E_2 = -1
    \end{equation*}
    Then:
    \begin{equation*}
        E = 0
    \end{equation*}
    They cancel each other. That is interference. So the fundamental thing that interference is \textbf{the amplitude}: $E_1 + E_2$.
\end{deepeningbox}

\highspace
Suppose light hits a screen. We don't see the electric field $E$ directly:
\begin{equation*}
    E = +3, \quad E = -2, \quad E = +1, \quad \dots
\end{equation*}
Because the electromagnetic field oscillates extremely rapidly, and our eyes cannot follow it. Instead, the screen receives \textbf{energy} from the light.
\begin{itemize}
    \item A region receiving more energy appears brighter.
    \item A region receiving less energy appears darker.
\end{itemize}
The quantity describing this is called the \definition{Intensity} $I$, the \textbf{energy delivered per unit area per unit time}, roughly speaking.

\highspace
For a light wave, intensity is proportional to $\propto$ the square of the wave amplitude:
\begin{equation*}
    I \propto \left|E\right|^2
\end{equation*}
\textcolor{Green3}{\faIcon{question-circle} \textbf{Why square the amplitude?}} One immediate reason the square makes sense is that amplitude can have a sign (e.g., $E = +1$ or $E = -1$), but \textbf{energy cannot be negative}. Squaring guarantees that the intensity is always positive. The sign is important for \textbf{interference}, but the final measured energy/intensity is always positive. Thus:
\begin{enumerate}
    \item Add the amplitudes of the waves to get the total amplitude.
    \item Square the total amplitude to get the intensity.
\end{enumerate}

\highspace
\begin{examplebox}[: Interference of two waves]
    Imagine each wave alone would have amplitude magnitude 1. If only wave 1 existed:
    \begin{equation*}
        E_1 = 1
    \end{equation*}
    So its intensity would be proportional to:
    \begin{equation*}
        I_1 \propto 1^{2} = 1
    \end{equation*}
    Similarly, if only wave 2 existed:
    \begin{equation*}
        E_2 = 1
    \end{equation*}
    So its intensity would be proportional to:
    \begin{equation*}
        I_2 \propto 1^{2} = 1
    \end{equation*}
    We might naïvely think that two waves together should therefore give:
    \begin{equation*}
        I = I_1 + I_2 = 1 + 1 = 2
    \end{equation*}
    But that is \textbf{not necessarily true}. We must first add their amplitudes.
    \begin{itemize}
        \item \textbf{Case A: the waves reinforce each other}. If:
            \begin{equation*}
                E_1 = 1, \quad E_2 = 1
            \end{equation*}
            Then:
            \begin{equation*}
                E = E_1 + E_2 = 1 + 1 = 2
            \end{equation*}
            So the intensity is:
            \begin{equation*}
                I \propto \left|E\right|^{2} = 2^{2} = 4
            \end{equation*}
            We have \textbf{constructive interference}.

            And if we had:
            \begin{equation*}
                E_1 = -1, \quad E_2 = -1
            \end{equation*}
            Then:
            \begin{equation*}
                E = E_1 + E_2 = -1 - 1 = -2
            \end{equation*}
            So the intensity is:
            \begin{equation*}
                I \propto \left|E\right|^{2} = \left|-2\right|^{2} = 2^{2} = 4
            \end{equation*}
            Again, we have \textbf{constructive interference}, even though the amplitudes are negative. The sign of the amplitude does not matter for the intensity.

        \item \textbf{Case B: the waves cancel each other}. If:
            \begin{equation*}
                E_1 = 1, \quad E_2 = -1
            \end{equation*}
            Then:
            \begin{equation*}
                E = E_1 + E_2 = 1 - 1 = 0
            \end{equation*}
            So the intensity is:
            \begin{equation*}
                I \propto \left|E\right|^{2} = 0^{2} = 0
            \end{equation*}
            So the screen is dark there. This is \textbf{destructive interference}.
    \end{itemize}
\end{examplebox}

\highspace
Let's generalize the previous example. Imagine one sinusoidal wave reaching some point $P$:
\begin{equation*}
    E_1(t) = E_0 \cos(\omega t)
\end{equation*}
Its average intensity is proportional to the square of the amplitude:
\begin{equation*}
    I_1 \propto E_0^2
\end{equation*}
So far, nothing strange. Now suppose \textbf{two waves} reach the same point $P$. Because ordinary wave equations are linear, the physical field is the sum:
\begin{equation*}
    E(t) = E_1(t) + E_2(t)
\end{equation*}
This is the classical \textbf{superposition principle for waves}. But intensity is not proportional to $E$, but rather to the \textbf{square} of the resulting amplitude:
\begin{equation*}
    I \propto \left|E\right|^{2} \implies I \propto \left|E_1 + E_2\right|^{2}
\end{equation*}
Expanding:
\begin{equation*}
    \left|E_1 + E_2\right|^{2} = \left|E_1\right|^{2} + \left|E_2\right|^{2} + 2\Re\left(E_1^* E_2\right)
\end{equation*}
The last term $2\Re\left(E_1^* E_2\right)$ is called the \definition{Inference Term}. This term is the reason why:
\begin{equation*}
    \left|\psi_1 + \psi_2\right|^{2} \neq \left|\psi_1\right|^{2} + \left|\psi_2\right|^{2}
\end{equation*}
Or in words, the \hl{intensity of two waves together is not equal to the sum of their individual intensities.} The inference term can be:
\begin{itemize}
    \item \definition{Constructive Interference} if the phases agree:
        \begin{equation*}
            2\Re\left(E_1^* E_2\right) > 0
        \end{equation*}
        The two contributions reinforce each other. We have a \textbf{maximum} in the intensity.

    \item \definition{Destructive Interference} if the phases disagree:
        \begin{equation*}
            2\Re\left(E_1^* E_2\right) < 0
        \end{equation*}
        The contributions oppose each other and may cancel. We have a \textbf{minimum} in the intensity.
\end{itemize}
So an \hl{alternating pattern of maxima and minima is a signature of interference}, and it is a \hl{signature of wave-like behavior.} Therefore, observing such a diffraction/interference pattern provides an experimental criterion for identifying the wave-like nature of light.

\highspace
\begin{deepeningbox}[: The interference term]
    Let's remove for a moment the absolute value from each term. Then the intensity is still proportional to the square of the amplitude:
    \begin{equation*}
        I \propto \left(E_1 + E_2\right)^{2}
    \end{equation*}
    Expanding it exactly like ordinary algebra:
    \begin{equation*}
        \left(E_1 + E_2\right)^{2} = E_1^{2} + E_2^{2} + 2 E_1 E_2
    \end{equation*}
    The last term $2 E_1 E_2$ is the \definition{Interference Term} in the simplest real-valued case. It is really \textbf{important} because without it, we would just have:
    \begin{equation*}
        E_1^{2} + E_2^{2}
    \end{equation*}
    Which is just ``intensity from wave 1 plus intensity from wave 2''. The extra term tells us what happens when the two waves meet.

    \highspace
    For example, if $E_1 = 1$ and $E_2 = 1$, then the intensity is:
    \begin{equation*}
        I \propto \left(1 + 1\right)^{2} = 4
    \end{equation*}
    But separately, the intensity from each wave is:
    \begin{equation*}
        I_1 + I_2 = 1^{2} + 1^{2} = 2
    \end{equation*}
    The extra term $+2$ comes from the interference term:
    \begin{equation*}
        2 E_1 E_2 = 2 \cdot 1 \cdot 1 = 2
    \end{equation*}
    So the waves reinforce each other. Now take $E_1 = 1$ and $E_2 = -1$. Then the intensity is:
    \begin{equation*}
        I \propto \left(1 - 1\right)^{2} = 0
    \end{equation*}
    But separately, the intensity from each wave is:
    \begin{equation*}
        I_1 + I_2 = 1^{2} + (-1)^{2} = 1 + 1 = 2
    \end{equation*}
    The extra term $-2$ comes from the interference term:
    \begin{equation*}
        2 E_1 E_2 = 2 \cdot 1 \cdot (-1) = -2
    \end{equation*}
    So the interference term is simply what can either increase the intensity or decrease it.

    \highspace
    \textcolor{Green3}{\faIcon{question-circle}} Now, why did we previously write the interference term as $2\Re\left(E_1^* E_2\right)$ instead of $2 E_1 E_2$? Because for light waves we often represent the amplitude using complex numbers. Then the correct version of the same algebra becomes:
    \begin{equation*}
        \begin{aligned}
            \left|E\right|^{2} &= \left|E_1 + E_2\right|^{2} \\[.3em]
            &\downarrow \left|z\right|^{2} = z \cdot z^* = \left(x + iy\right)\left(x - iy\right) = x^2 + y^2 \\[.3em]
            &\downarrow \text{ if we assume } E_1 + E_2 = z \\[.3em]
            &= \left(E_1 + E_2\right)\left(E_1 + E_2\right)^* \\[.3em]
            &= \left(E_1 + E_2\right)\left(E_1^* + E_2^*\right) \\[.3em]
            &= E_1 E_1^* + E_1 E_2^* + E_2 E_1^* + E_2 E_2^* \\[.3em]
            &\downarrow E_1 E_1^* = \left|E_1\right|^{2}, \quad E_2 E_2^* = \left|E_2\right|^{2} \\[.3em]
            &= \left|E_1\right|^{2} + \left|E_2\right|^{2} + E_1 E_2^* + E_2 E_1^*
        \end{aligned}
    \end{equation*}
    Now the last two terms are complex conjugates of each other. Let:
    \begin{equation*}
        z = E_1 E_2^*
    \end{equation*}
    Thus its complex conjugate is:
    \begin{equation*}
        z^* = \left(E_1 E_2^*\right)^* = E_1^* \left(E_2^*\right)^* = E_1^* E_2
    \end{equation*}
    Therefore:
    \begin{equation*}
        E_1 E_2^* + E_2 E_1^* = z + z^*
    \end{equation*}
    And for any complex number:
    \begin{equation*}
        z = a + ib
    \end{equation*}
    We have:
    \begin{equation*}
        z^* = a - ib
    \end{equation*}
    So:
    \begin{equation*}
        z + z^* = a + ib + a - ib = 2a
    \end{equation*}
    But $a$ is the real part of $z$, represented as $\Re(z)$. So we have:
    \begin{equation*}
        z + z^* = 2\Re(z)
    \end{equation*}
    Therefore:
    \begin{equation*}
        E_1 E_2^* + E_2 E_1^* = 2\Re\left(E_1 E_2^*\right) \equiv 2\Re\left(E_1^* E_2\right)
    \end{equation*}
    And conceptually it is just the same cross term as $2 E_1 E_2$.
\end{deepeningbox}

\begin{figure}[!htp]
    \centering
    \tikzsetnextfilename{diffraction-pattern}
    \tikzwidth{\linewidth}
    \begin{tikzpicture}[
            >=Stealth,
            wave/.style={blue!60!black, line width=0.5pt},
            ray/.style={red!60!black, line width=0.4pt},
            every node/.style={font=\small}
        ]

        %% ------------------------------------------------------------
        %% Geometry: plane wave -> two slits -> screen
        %% ------------------------------------------------------------
        \def\slit{0.7}      % half-distance between the two slits
        \def\Lscr{6}        % barrier-to-screen distance
        \def\Ymax{3}        % half-height of the screen

        % incoming plane wave fronts
        \foreach \x in {-2.4,-1.8,-1.2,-0.6}
        \draw[wave] (\x,-2.2) -- (\x,2.2);
        \draw[->, thick] (-2.9,0) -- (-2.5,0);
        \node[above] at (-1.5,2.3) {incident plane wave};

        % barrier with two slits (gaps of width 0.12 at y = +-\slit)
        \fill[gray!70] (-0.06,-\Ymax) rectangle (0.06,-\slit-0.06);
        \fill[gray!70] (-0.06,-\slit+0.06) rectangle (0.06,\slit-0.06);
        \fill[gray!70] (-0.06,\slit+0.06) rectangle (0.06,\Ymax);
        \node[above] at (0,\Ymax) {two slits};

        % circular wavefronts emitted by each slit
        \begin{scope}
            \clip (0.1,-\Ymax) rectangle (\Lscr-0.1,\Ymax);
            \foreach \r in {0.9,1.6,2.3,3.0,3.7,4.4,5.1}{
                \draw[wave] (0,\slit)  ++(-60:\r) arc[start angle=-60,end angle=60,radius=\r];
                \draw[wave] (0,-\slit) ++(-60:\r) arc[start angle=-60,end angle=60,radius=\r];
            }
        \end{scope}

        % a few rays towards a maximum (y=0) and towards a generic point
        \draw[ray] (0,\slit)  -- (\Lscr,0);
        \draw[ray] (0,-\slit) -- (\Lscr,0);
        \draw[ray] (0,\slit)  -- (\Lscr,1.8);
        \draw[ray] (0,-\slit) -- (\Lscr,1.8);

        % screen with grey-scale fringes: darkness = 1 - normalised intensity
        % I(y) = sinc^2(a y) * cos^2(b y)
        \foreach \k in {-60,...,59}{
            \pgfmathsetmacro{\y}{0.05*\k}
            \pgfmathsetmacro{\yy}{\y+0.025}
            \pgfmathsetmacro{\I}{(sin(deg(1.2*\yy))/(1.2*\yy))^2 * (cos(deg(4*\yy)))^2}
            \pgfmathtruncatemacro{\g}{round(100*(1-\I))}
            \fill[black!\g] (\Lscr,\y) rectangle (\Lscr+0.3,\y+0.05);
        }
        \draw (\Lscr,-\Ymax) rectangle (\Lscr+0.3,\Ymax);
        \node[below] at (\Lscr+0.15,-\Ymax) {screen};

        % intensity profile I(y) drawn to the right of the screen
        \def\x0{\Lscr+0.7}
        \draw[->] (\x0,-\Ymax) -- (\x0,\Ymax+0.3) node[above] {$y$};
        \draw[->] (\x0,0) -- (\x0+2.3,0) node[right] {$I$};
        \draw[thick, red!70!black, samples=400, domain=-\Ymax:\Ymax, smooth, variable=\y]
        plot ({\x0 + 2*(sin(deg(1.2*\y))/(1.2*\y))^2 * (cos(deg(4*\y)))^2}, {\y});
        % single-slit envelope
        \draw[dashed, gray, samples=200, domain=-\Ymax:\Ymax, smooth, variable=\y]
        plot ({\x0 + 2*(sin(deg(1.2*\y))/(1.2*\y))^2}, {\y});
        \node[gray, anchor=west] at (\x0+1.4,2.2) {envelope};

    \end{tikzpicture}
    \caption{\textbf{Diffraction and interference in a two-slit experiment.} An incident plane wave passes through two slits; the emerging waves spread, overlap, and interfere, producing an intensity pattern on the screen modulated by a diffraction envelope.}
    \label{fig:diffraction-pattern}
\end{figure}

\highspace
\begin{deepeningbox}
    Some details about the \autoref{fig:diffraction-pattern}:
    \begin{itemize}
        \item[\textcolor{Green3}{\faIcon{question-circle}}] \textcolor{Green3}{\textbf{What is the ``incident plane wave''?}} On the left we have several \textbf{parallel blue vertical lines}, those lines are \textbf{wavefronts of the same incoming wave}.

            A wavefront is a set of points where the wave is in the same state of oscillation. We do not need to know particular details; visually, think of each blue line as one crest/front of the wave. Because these wavefronts are straight and parallel, we call it a \textbf{plane wave}. So ``\emph{incident plane wave}'' simply means \hl{a wave arriving at the barrier with flat, parallel wavefronts.} ``\emph{Incident}'' just means \textbf{incoming}.

            Finally, the arrow indicates that the wave propagates from left to right.

        \item[\textcolor{Green3}{\faIcon{question-circle}}] \textcolor{Green3}{\textbf{What are the ``two slits''?}} The grey vertical object is a barrier with \textbf{two small openings}. Most of the wave is blocked by the barrier, only the portions reaching those openings pass through.

            Here, the crucial observation is what happens \textbf{after} each slit. Instead of continuing only straight ahead, the wave \textbf{spreads out}. That's why the figure draws curved blue wavefronts coming from each slit.

            This spreading is the \textbf{diffraction} part! Each slit effectively gives us a new spreading wave. And now we have precisely what we need fo \textbf{interference}: two waves reaching the same points on the screen (black rectangle on the right).

        \item[\textcolor{Green3}{\faIcon{question-circle}}] \textcolor{Green3}{\textbf{What happens between the slits and the screen?}} Pick some point $P$ on the screen. There are two wave contributions reaching $P$: $\psi_1$ from the upper slit and $\psi_2$ from the lower slit. The red straight lines in the figure illustrate examples of those two paths.

            The two contributions add:
            \begin{equation*}
                \psi(P) = \psi_1(P) + \psi_2(P)
            \end{equation*}
            But depending on the point $P$, the two paths have different lengths. So the waves may arrive ``\emph{in phase}'' (reinforce, constructive interference), ``\emph{out of phase}'' (cancel out, destructive interference), or somewhere in between (weaker reinforcement or weaker cancellation). And that's why different points of the screen have different brightness.

        \item[\textcolor{Green3}{\faIcon{question-circle}}] \textcolor{Green3}{\textbf{Why is the center of the screen bright?}} Call the two slit-to-screen distances $r_1$ and $r_2$. At the exact center of the screen, by symmetry, the two distances are equal:
            \begin{equation*}
                r_1 = r_2
            \end{equation*}
            Therefore the two waves have traveled the same distance, and they reinforce each other:
            \begin{equation*}
                \psi_1 + \psi_2 \implies I \propto \left|\psi_1 + \psi_2\right|^{2}
            \end{equation*}
            So the intensity is large and the \hl{central bright band corresponds to constructive interference.}

        \item[\textcolor{Green3}{\faIcon{question-circle}}] \textcolor{Green3}{\textbf{What are the dark bands?}} Move upward or downward along the screen. Now one slit is slightly closer than the other:
            \begin{equation*}
                r_1 \neq r_2
            \end{equation*}
            At certain position, the difference in traveled distance causes one wave to arrive exactly opposite in phase to the other.

            For example:
            \begin{equation*}
                \psi_1 = A, \qquad \psi_2 = -A
            \end{equation*}
            Therefore:
            \begin{equation*}
                \psi_1 + \psi_2 = A - A = 0
            \end{equation*}
            And hence the intensity is zero:
            \begin{equation*}
                I \propto \left|\psi_1 + \psi_2\right|^{2} = 0
            \end{equation*}
            So that location is dark, corresponding to \hl{destructive interference.}

        \item[\textcolor{Green3}{\faIcon{question-circle}}] \textcolor{Green3}{\textbf{What is that red curve on the right?}} The horizontal direction of the red graph represents the \textbf{intensity} $I$ measured on the screen, while vertically we move along the screen coordinate $y$.

            So where the solid red curve extends far to the right, the \hl{intensity is large and the screen is bright.} Where it comes \hl{close to zero $I \approx 0$, the screen is dark.}

            So the oscillating solid red curve is essentially a \hl{graph of $I(y)$, the intensity as a function of vertical position on the screen.}

        \item[\textcolor{Green3}{\faIcon{question-circle}}] \textcolor{Green3}{\textbf{What is the ``envelope''?}} Look at the \textbf{dashed red curve} labelled ``envelope''. It is \textbf{not} ``\emph{the level of interference}''. Instead, it \hl{describes how large the interference peaks are allowed to become at different positions.}

            A more natural question is: \textbf{\emph{why are the bright fringes near the center stronger than the bright fringes farther away?}} What is exactly what the dashed curve in figure is trying to show. The red solid curve $I(y)$ oscillates between bright and dark because of the \textbf{interference between the two slits}.

            But notice also that the heights of the bright peaks are not all the same. Near the center, the peaks are large; farther away, they become smaller. The dashed curve simply joins the tops of those possible peaks. So the \textbf{envelope} is the \hl{smooth curve giving the maximum possible intensity at each position on the screen (at each $y$).} In other words, it is just a boundary for the interference pattern.

            \textcolor{Green3}{\faIcon{question-circle} \textbf{But why does the envelope exist?}} Because the slits are not mathematical points; they have a finite width. Even \textbf{one slit by itself} produces a diffraction (i.e., spreading) pattern on the screen. So each slit does not send the same wave strength equally in every direction. it sends most strongly approximately forward, around the center of the screen, and more weakly at larger angles.

            Therefore, before we even combine the two slits, each slit already has something like: strong near the center, weaker farther away. That broad distribution is the \textbf{single-slit diffraction pattern}. Now, if we combine the two slits, the two effects happen simultaneously and each slit diffracts; then the waves from the two slits interfere and the final pattern is the product of the two effects:
            \begin{itemize}
                \item \textbf{Single-slit diffraction} gives the broad envelope.
                \item \textbf{Two-slit interference} gives the oscillating pattern within that envelope.
            \end{itemize}
    \end{itemize}
\end{deepeningbox}

\begin{figure}[!htp]
    \centering
    \tikzsetnextfilename{interference-term}
    \begin{tikzpicture}[
            >=Stealth,
            every node/.style={font=\small},
            x=1.5cm, y=1.6cm
        ]

        %% ------------------------------------------------------------
        %% The interference term: |psi1+psi2|^2 versus |psi1|^2+|psi2|^2
        %% ------------------------------------------------------------
        \draw[->] (-3.3,0) -- (3.5,0) node[right] {$\phi_1-\phi_2$};
        \draw[->] (0,-0.2) -- (0,2.3) node[right] {$I/A^2$};
        \foreach \n/\lab in {-3/-3\pi,-2/-2\pi,-1/-\pi,1/\pi,2/2\pi,3/3\pi}
        \draw (\n,0.05) -- (\n,-0.05) node[below] {$\lab$};
        \node[below left] at (0,0) {$0$};
        \draw (-0.05,2) -- (0.05,2) node[left=4pt] {$4$};
        \draw (-0.05,1) -- (0.05,1) node[left=4pt] {$2$};

        % with interference: 2A^2 (1 + cos(delta))  ->  in units of A^2 divided by 2 for the plot
        \draw[thick, red!70!black, samples=300, domain=-3:3, smooth, variable=\x]
        plot ({\x}, {1+cos(180*\x)});
        % without interference: |psi1|^2 + |psi2|^2 = 2 A^2
        \draw[thick, dashed, blue!60!black] (-3.2,1) -- (3.2,1);

        \node[red!70!black, above] at (0,2) {constructive};
        \node[red!70!black, above] at (2,2) {constructive};
        \node[red!70!black, above] at (-2,2) {constructive};
        \node[red!70!black, below] at (1,0.02) [yshift=-0.55cm] {destructive};
        \node[red!70!black, below] at (-1,0.02) [yshift=-0.55cm] {destructive};

        \node[blue!60!black, anchor=west] at (3.2,1.25)
        {$|\psi_1|^2+|\psi_2|^2$};
        \node[red!70!black, anchor=west] at (3.2,0.55)
        {$|\psi_1+\psi_2|^2$};

    \end{tikzpicture}
    \caption{\textbf{The interference term.} The resulting intensity depends on the relative phase $\phi_1-\phi_2$: constructive interference produces maxima, while destructive interference produces minima.}
\end{figure}
